El presente capítulo tiene como finalidad establecer las políticas, procedimientos y responsabilidades que rigen la disposición final, almacenamiento temporal, manejo y eliminación segura de los componentes, repuestos y baterías utilizadas en los Sistemas de Aeronaves No Tripuladas (UAS) de SKYFILMS.

Este procedimiento busca garantizar el cumplimiento de las disposiciones establecidas por la Unidad Administrativa Especial de Aeronáutica Civil (Aerocivil), y la reglamentación ambiental vigente, asegurando la prevención de riesgos eléctricos, incendios, contaminación o daño ambiental.

\section{Alcance}
Este procedimiento aplica a todos los componentes reemplazados, baterías de polímero de litio (LiPo), accesorios electrónicos, sensores, estructuras dañadas, y elementos consumibles generados durante las actividades de mantenimiento preventivo y correctivo de los UAS operados por SKYFILMS.

Comprende desde el retiro del componente o batería, su identificación, almacenamiento temporal, evaluación técnica y clasificación, hasta su entrega a un gestor autorizado o eliminación definitiva.

\section{Política General}
SKYFILMS adopta como principio técnico y ambiental que todo componente o batería retirada de servicio debe ser:

\begin{enumerate}
    \item Identificada y registrada adecuadamente antes de su disposición final.

    \item Almacenada en condiciones seguras, evitando riesgos eléctricos, químicos o de incendio.

    \item Evaluada técnicamente para determinar si puede ser reutilizada, reparada o debe ser desechada.

    \item Eliminada o reciclada únicamente a través de entidades o gestores ambientales autorizados, conforme a las regulaciones del Ministerio de Ambiente y Desarrollo Sostenible y la Aerocivil.
\end{enumerate}

El cumplimiento de esta política es obligatorio para todo el personal técnico, administrativo y contratistas que intervengan en los procesos de mantenimiento o eliminación de materiales.

\section{Procedimiento General de Disposición}

\begin{enumerate}
    \item \textbf{Identificación y retiro del componente o batería:} Una vez retirado de la aeronave, el técnico debe etiquetar el elemento con la información básica: fecha, tipo, número de serie, motivo del retiro y estado \textbf{(dañado, vencido, deteriorado)}.

    \item \textbf{Registro y clasificación:} El componente se registra en el \textbf{Formato de Control de Disposición Final}, donde se indica su condición técnica (reparable, reutilizable o no aeronavegable).

    \item \textbf{Evaluación técnica:} El \textbf{Tecnico UAS} evalúa si el componente puede ser reutilizado o si debe ser descartado definitivamente.

    \item \textbf{Almacenamiento temporal:} Los componentes y baterías retiradas se almacenan en el \textbf{Área de Material Retirado}, en estanterías o contenedores identificados, aislados del área de mantenimiento activo.

    \item \textbf{Disposición final o reciclaje:} Los elementos declarados como \textbf{no aeronavegables o no reutilizables} se entregan a gestores ambientales o proveedores certificados para su eliminación o reciclaje conforme a la normatividad ambiental. La entrega se documenta mediante acta o certificado de disposición, archivada junto con los registros de mantenimiento.

    \item \textbf{Actualización de registros:} Todo proceso de eliminación o entrega debe quedar registrado en el \textbf{Anexo 5 - Disposición Final de Componentes}, controlado por el \textbf{Tecnico UAS}.
\end{enumerate}

\section{Manejo y Almacenamiento de Baterías de Litio}
Las \textbf{baterías de polímero de litio (LiPo)} utilizadas por SKYFILMS requieren tratamiento especial debido a su potencial riesgo eléctrico y químico. Para su manipulación y eliminación se cumplirán las siguientes disposiciones:

\begin{itemize}
    \item Las baterías dañadas, hinchadas o con fugas deberán retirarse de inmediato del servicio y aislarse en bolsas o mantas ignífugas dentro de un recipiente cerrado.

    \item Las áreas de almacenamiento temporal deberán estar ventiladas, señalizadas y equipadas con un extintor ABC y mantas ignífugas.

    \item El nivel de carga de las baterías almacenadas no debe exceder el 30\% de su capacidad total para reducir riesgos de reacción térmica.

    \item Está prohibido depositar baterías en basureros comunes o mezclarlas con residuos no peligrosos.

    \item Las baterías descartadas se entregarán a gestores autorizados de residuos peligrosos, de acuerdo con el Decreto 1076 de 2015 del Ministerio de Ambiente.
\end{itemize}

\section{Disposición Final de Componentes Electrónicos y Mecánicos}
Los componentes electrónicos, motores, hélices, tarjetas de control, sensores, estructuras y demás piezas retiradas se gestionarán conforme a su tipo:

\begin{itemize}
    \item \textbf{Componentes recuperables:} se almacenarán para análisis y posible reparación, bajo registro técnico de control.

    \item \textbf{Componentes no aeronavegables o irrecuperables:} se clasificarán como residuos electrónicos (RAEE) y se entregarán a un gestor certificado.

    \item \textbf{Materiales plásticos o metálicos} podrán ser enviados a reciclaje si cumplen condiciones seguras de manejo y no presentan contaminación química o eléctrica.
\end{itemize}

Cada disposición deberá contar con soporte documental que permita la trazabilidad completa desde el retiro hasta su eliminación final.

\section{Control Documental y Registros}
Los siguientes registros respaldan el cumplimiento de este procedimiento:

\begin{itemize}
    \item Formato de Control de Disposición Final de Componentes y Baterías.

    \item Certificados de Entrega a Gestores Autorizados.

    \item Registro fotográfico o documental del estado del elemento al momento del retiro (cuando aplique).
\end{itemize}

Todos los registros se conservan bajo custodia del \textbf{Técnico de UAS}, y estarán disponibles para auditorías internas o inspecciones de la Aerocivil.

\section{Responsabilidades}
\begin{itemize}
    \item \textbf{Gerente General:} asegura el cumplimiento de la normativa ambiental y la disponibilidad de recursos para la gestión adecuada de residuos técnicos.

    \item \textbf{Técnico de UAS:} 
    
        \begin{itemize}
            \item Supervisa la correcta aplicación del procedimiento y aprueba las disposiciones finales de materiales.

            \item Lleva el control documental, coordina la entrega a gestores y mantiene los registros actualizados.
            
            \item Ejecuta la segregación, identificación y registro inicial de los elementos retirados.
        \end{itemize}
 
\end{itemize}

\section{Disposiciones Finales}
El manejo responsable de componentes y baterías retiradas es una obligación técnica y ambiental de SKYFILMS, y su cumplimiento es condición indispensable para mantener la conformidad con los Reglamentos Aeronáuticos de Colombia (RAC 100) y las políticas del Sistema de Gestión de Seguridad Operacional (SMS).

El incumplimiento de los procedimientos establecidos en este capítulo será considerado una no conformidad técnica grave, y estará sujeto a acciones correctivas inmediatas y posibles sanciones internas.

El proceso descrito garantiza que SKYFILMS opere bajo los principios de seguridad, sostenibilidad ambiental y cumplimiento normativo, manteniendo el compromiso institucional con la Aerocivil y con la preservación del entorno.